\documentclass{article}
\usepackage[utf8]{inputenc}
\usepackage[a4paper, total={6in, 8in}]{geometry}
\usepackage{graphicx}
\usepackage{amsmath}
\usepackage{enumitem}



\title{Econ 20000-Problem Set 3}
\author{Dr. Fulya Ersoy}
\date{}
\begin{document}
\maketitle
\begin{enumerate}
\item (Quasilinear utility Function-Marshallian Demand) Lima 5.8 

\includegraphics[width=\textwidth]{Lima-Exercise-5.8.jpg}

\item[] a. $\max_{x,y}U(x,y)=y+\sqrt{x}$, w.r.t budget constraint: $p_xx+p_yy=m$.
\item[] b. $\mathcal{L}=y+\sqrt{x}+\lambda(m-p_xx-p_yy)$. First order conditions:
\item[] $\frac{\partial \mathcal{L}}{\partial x}=\frac{1}{2\sqrt x}-\lambda p_x=0$. $\frac{\partial \mathcal{L}}{\partial y}=1-\lambda p_y=0$. $\frac{\partial \mathcal{L}}{\partial \lambda}=m-p_xx-p_yy=0$.
\item[] c. From first two FOCs: $\lambda=\frac{1}{p_y}=\frac{1}{2\sqrt{x}p_x} \Rightarrow MRS_{xy}=\frac{MU_x}{MU_y}=\frac{p_x}{p_y}=\frac{1}{2\sqrt{x}}$. Thus, consumer indifference curves decrease with $x$, in other words, ICs are convex and quasilinear in y.
\item[] d. From MRS condition: $\frac{1}{2\sqrt{x}}=\frac{p_x}{p_y}\Rightarrow x^*=(\frac{p_y}{2p_x})^2$. Substituting into budget constraint, $p_xx^*+p_yy^*=m \Rightarrow y^*=\frac{m-p_xx^*}{p_y}=\frac{m-p_x(\frac{p_y}{2p_x})^2}{p_y}=\frac{m}{p_y}-\frac{p_y}{4p_x}$.
\item[] e. Substituting given values: $x^*=(\frac{1}{2(1)})^2=\frac{1}{4},\ y^*=\frac{\frac{1}{9}}{1}-\frac{1}{4(1)}=-\frac{5}{36}$. Since $y^*$ cannot be negative by the feasibility condition, the consumer spends all income on the good with higher relative utility (corner vs interior solution, $y=0$). Here $U(x,0)=\sqrt{x}, $ subject to $x=\frac{m}{p_x}=\frac{1}{9}$. Thus, optimal consumption levels are $x=\frac{1}{9},\ y=0$ with $U=\frac{1}{3}$.
\item[] f. If income increases to $m=1$, interior solution works: $x^*=(\frac{1}{2(1)})^2=\frac{1}{4}$, $y^*=\frac{1}{1}-\frac{1}{4(1)}=\frac{3}{4}$.
\item[] g. $x^*$ is not affected by changes in income. Share of income spent on good $x$ falls with greater income because $x$ is fixed at $p_xx^*$. $y^*$ increases linearly with income. Share of income spent on good $y$ rises with greater income, inverse to relationship with $x$.
\end{enumerate}
\begin{enumerate}

 \setcounter{enumi}{1}
\item (Perfect Subtitutes -Marshallian Demand) Suppose that indifference curves are described by straight lines with a slope of -b. Given prices ($p_1$ and $p_2$) and income ($m$), find the consumers demand function for good 1 and good 2. Hint: This is a case where the Lagrangian method doesn't work.
\item[] A utility with $-b$ slope is $u(x,y)=bx+y$, maximized by the budget $p_1x+p_2y\leq m$. Here, $MU_1=b,\ MU_2=1$ and $\frac{MU_1}{p_1}=\frac{b}{p_1},\ \frac{MU_2}{p_2}=\frac{1}{p_2}$. There are three cases in the demand function:
\item[] If $\frac{b}{p_1}>\frac{1}{p_2}$, consumer only buys good 1: $x^*_1(p_1,p_2,m)=\frac{m}{p_1},\ x^*_2=0$.
\item[] If $\frac{b}{p_1}<\frac{1}{p_2}$, consumer only buys good 2: $x^*_1=0,\ x^*_2(p_1,p_2,m)=\frac{m}{p_2}$.
\item[] If $\frac{b}{p_1}=\frac{1}{p_2}$, every bundle on the budget line is optiomal: $\{(x,y)\geq0,\ p_1x+p_2y= m\}$.
\item Consider the problem:
$\underset{x, y}{min} (p_x x+p_y y)$ subject to $\alpha lnx+ (1-\alpha) lny=ln\bar{u}$
\begin{enumerate}[label=(\alph*)]
\item Form the Lagrangian function associated with this constrained minimization problem.
\item[] Choose multiplier $\eta$. $\mathcal{L}(x,y,\eta)=p_xx+p_yy+\eta(ln\bar u-\alpha lnx-(1-\alpha)lny=0)$
\item Take the partial derivatives with respect to x, y, and the lagrange multiplier ($\eta$).
\item[] $\frac{\partial \mathcal{L}}{\partial x}=p_x-\eta\frac{\alpha}{x}$, $\frac{\partial \mathcal{L}}{\partial y}=p_y-\eta\frac{(1-\alpha)}{y}$, $\frac{\partial \mathcal{L}}{\partial \eta}=ln\bar u-\alpha lnx-(1-\alpha)lny$.
\item Equate the partial derivatives to zero to find the first order conditions and solve the system of equations. Find the Hicksian demand functions and label them $x_h^*$ and $y_h^*$.
\item[] Set first 2 derivatives $=0$: $p_x=\eta\frac{\alpha}{x}\Rightarrow \eta=\frac{p_xx}{\alpha}$, $p_y=\eta\frac{1-\alpha}{y}\Rightarrow \eta=\frac{p_yy}{1-\alpha}$.
\item[] Equate expressions for $\eta$: $\frac{p_xx}{\alpha}=\frac{p_yy}{1-\alpha} \Rightarrow \frac{x}{y}=\frac{\alpha p_y}{(1-\alpha)p_x}$. Set arb $r=\frac{\alpha p_y}{(1-\alpha)p_x}\Rightarrow x=ry$.
\item[] Sub into constraint for $y$: $x^\alpha+y^{1-\alpha}=\bar u \Rightarrow (ry)^\alpha+y^{1-\alpha}=\bar u \Rightarrow r^\alpha y=\bar u \Rightarrow y=\bar ur^{-\alpha}$.
\item[] Thus, $y_h^*(p_x,p_y,\bar u)=\bar u(\frac{\alpha p_y}{(1-\alpha)p_x})^{-\alpha}=\bar u(\frac{(1-\alpha)p_x}{\alpha p_y})^{\alpha}$.
\item[] And, $x^*_h(p_x,p_y,\bar u)=r\cdot y^*_h=r\cdot\bar ur^{-\alpha}=\bar u(\frac{\alpha p_y}{(1-\alpha)p_x})^{1-\alpha}$.
\item Form the expenditure function. Write a single sentence to describe what this function tells us. Use Shephard’s lemma to verify that the Hicksian demand functions are the same as what you found in the previous part.
\item[] $e=p_xx_h^*+p_yy_h^* \Rightarrow e=\bar up^\alpha_xp_y^{1-\alpha}(\frac{\alpha^{1-\alpha}}{(1-\alpha)^{1-\alpha}}+\frac{(1-\alpha)^{\alpha}}{\alpha^{\alpha}}) \Rightarrow e(p_x,p_y,\bar u)=\bar u(\frac{p_x^\alpha p_y^{1-\alpha}}{\alpha^\alpha(1-\alpha)^{1-\alpha}})$.
\item[] The known expenditure function depicts the minimal amount of money needed to be spent at given market prices ($p_x,p_y$) to achieve a target utility $\bar u$.
\item[] Shephard's: $x_h^*=\frac{\partial e}{\partial p_x}=\bar u(\frac{\alpha p_x^{\alpha-1}p_y^{1-\alpha}}{\alpha^\alpha (1-\alpha)^{1-\alpha}})=\frac{\alpha e}{p_x}$, $y_h^*=\frac{\partial e}{\partial p_y}=\bar u(\frac{(1-\alpha) p_x^{\alpha}p_y^{-\alpha}}{\alpha^\alpha (1-\alpha)^{1-\alpha}})=\frac{(1-\alpha) e}{p_y}$. Plugging back in confirms Hicksian demand functions for $x$ and $y$ as shown above.
\item Definition: A function $f(x, y)$ is said to be homogeneous of degree k if $f(t · x, t · y) = t^k · f(x, y)$, where k is a nonnegative integer and $t > 0$ is a positive constant. Use this definition to show that the expenditure function is homogenous of degree 1 in prices. Notice that the statement you are trying to show involves prices and NOT the target utility level. Briefly explain what this property means in words.
\item[] Checking $e(tp_x,tp_y,\bar u)$: $\Rightarrow \bar u(\frac{(tp_x)^\alpha(tp_y)^{1-\alpha}}{\alpha^\alpha(1-\alpha)^{1-\alpha}})=t\cdot\bar u(\frac{(p_x)^\alpha(p_y)^{1-\alpha}}{\alpha^\alpha(1-\alpha)^{1-\alpha}})=t\cdot e(p_x,p_y,\bar u)$.
\item[] Thus, $e$ is homogeneous degree 1. Scaling prices by $t$ scales expenditure by same factor $t$.
\end{enumerate}
\item Compute the expenditure function for the perfect complements utility function $u(x_1, x_2) = min\{x_1, x_2\}$ when a consumer faces prices $p_1$ and $p_2$ and wants to achieve target utility $\bar{u}$ . Does your result make sense? Why or why not?
\item[] To achieve $\bar u$ with $u=$ min$\{x_1,x_2\}$, $x_1\geq \bar u$ and $x_2\geq \bar u$. Then, the cost-minimizing method is to set $x_1=x_2=\bar u$. Hence, $e(p_1,p_2,\bar u)=p_1\bar u +p_2\bar u=\bar u(p_1+p_2)$. This function makes sense due to fixed 1:1 ratio per unit of utility, resulting in total cost $=\ \bar u \ \cdot \ $sum of prices.
 \end{enumerate}
\end{document}
